% ============================================================
%  sections/chapter_03.tex
%  Section 3: Optimization Methods
% ============================================================

\section{Optimization Methods}
\label{sec:optimization}

The Tucker approximation problem~\eqref{eq:general_opt} is
\emph{jointly non-convex} in the factor matrices
$\{\mat{U}^{(n)}\}$ and the core tensor $\tensor{G}$.
We study three algorithmic families that address this non-convexity in
different ways: Alternating Least Squares (ALS) with its higher-order
refinement HOOI for the Frobenius objective, Iteratively Reweighted Least
Squares (IRLS) for the $\ell_1$ objective, and gradient-based methods
for the Huber objective.

% ------------------------------------------------------------
\subsection{Alternating Least Squares and HOOI}
\label{subsec:als}
% ------------------------------------------------------------

\paragraph{ALS principle.}
Although $\mathcal{L}_F$ is jointly non-convex in all factors simultaneously,
it is \emph{quadratic} --- and hence convex --- when all but one factor
matrix are held fixed.
ALS exploits this structure by cycling over the factors and solving each
resulting least-squares problem exactly.

\paragraph{Mode-$n$ subproblem.}
Fix all factor matrices except $\mat{U}^{(n)}$ and define
\begin{equation}
  \mat{Z}^{(n)}
  \;\coloneqq\;
  \mat{G}_{(n)}
  \!\left(
    \mat{U}^{(N)} \kron \cdots \kron \mat{U}^{(n+1)}
    \kron \mat{U}^{(n-1)} \kron \cdots \kron \mat{U}^{(1)}
  \right)^{\!\top}
  \;\in\; \R^{r_n \times \prod_{k\neq n} I_k}.
  \label{eq:Z_factor}
\end{equation}
The objective reduces to
\begin{equation}
  \min_{\mat{U}^{(n)}} \;
  \bigl\|\unfold{X}{n} - \mat{U}^{(n)} \mat{Z}^{(n)}\bigr\|_F^2,
  \label{eq:als_subproblem}
\end{equation}
whose closed-form solution is
\begin{equation}
  \mat{U}^{(n)}
  \;\leftarrow\;
  \unfold{X}{n} \mat{Z}^{(n)\top}
  \!\left(\mat{Z}^{(n)} \mat{Z}^{(n)\top}\right)^{-1}.
  \label{eq:als_update}
\end{equation}
When the columns of $\mat{U}^{(n)}$ are additionally constrained to be
orthonormal --- as in Tucker --- the optimal solution is the matrix of
leading $r_n$ left singular vectors of $\unfold{X}{n} \mat{Z}^{(n)\top}$,
making the matrix inverse unnecessary.

\paragraph{HOOI algorithm.}
Higher-Order Orthogonal Iteration (HOOI)~\cite{de2000best} applies ALS
with orthonormal constraints.
At each iteration, the contracted tensor
\begin{equation}
  \tensor{Y}^{(n)}
  \;=\; \tensor{X}
        \times_1 {\mat{U}^{(1)}}^\top
        \cdots
        \times_{n-1} {\mat{U}^{(n-1)}}^\top
        \times_{n+1} {\mat{U}^{(n+1)}}^\top
        \cdots
        \times_N {\mat{U}^{(N)}}^\top
  \label{eq:hooi_contracted}
\end{equation}
is formed (omitting mode $n$), and the update is
\begin{equation}
  \mat{U}^{(n)}
  \;\leftarrow\;
  \text{leading-}r_n\text{ left singular vectors of }
  \mat{Y}^{(n)}_{(n)}.
  \label{eq:hooi_update}
\end{equation}
After all factor matrices have converged, the core is recovered as
\begin{equation}
  \tensor{G}
  \;=\; \tensor{X}
        \times_1 {\mat{U}^{(1)}}^\top
        \times_2 {\mat{U}^{(2)}}^\top
        \cdots
        \times_N {\mat{U}^{(N)}}^\top.
  \label{eq:hooi_core}
\end{equation}

Algorithm~\ref{alg:hooi} summarizes the full procedure.

\begin{algorithm}[t]
  \caption{HOOI for Tucker decomposition (Frobenius norm)}
  \label{alg:hooi}
  \KwIn{Tensor $\tensor{X} \in \R^{I_1 \times \cdots \times I_N}$,
        target rank $(r_1, \ldots, r_N)$,
        tolerance $\varepsilon > 0$,
        max iterations $T$}
  \KwOut{Core $\tensor{G}$, factor matrices $\{\mat{U}^{(n)}\}$}

  \BlankLine
  \tcp{Initialization via HOSVD}
  \For{$n \leftarrow 1$ \KwTo $N$}{
    $\mat{U}^{(n)} \leftarrow$ leading-$r_n$ left singular vectors of
    $\unfold{X}{n}$\;
  }
  \BlankLine
  \tcp{Iterative refinement}
  \For{$t \leftarrow 1$ \KwTo $T$}{
    \For{$n \leftarrow 1$ \KwTo $N$}{
      $\tensor{Y} \leftarrow
        \tensor{X} \times_{k \neq n} {\mat{U}^{(k)}}^\top$\;
      $\mat{U}^{(n)} \leftarrow$
        leading-$r_n$ left singular vectors of $\unfold{Y}{n}$\;
    }
    $\tensor{G} \leftarrow
      \tensor{X} \times_1 {\mat{U}^{(1)}}^\top \cdots
                 \times_N {\mat{U}^{(N)}}^\top$\;
    $e_t \leftarrow
      \frobn{\tensor{X} - \tensor{G} \times_1 \mat{U}^{(1)}
             \cdots \times_N \mat{U}^{(N)}}$\;
    \lIf{$|e_t - e_{t-1}| < \varepsilon$}{\textbf{break}}
  }
  \Return $\tensor{G},\; \{\mat{U}^{(n)}\}$\;
\end{algorithm}

\paragraph{Convergence.}
Each HOOI iteration decreases the Frobenius reconstruction error
monotonically, because each mode-$n$ update is an exact minimizer of the
subproblem~\eqref{eq:als_subproblem} given the other factors.
The sequence $\{e_t\}$ is therefore non-increasing and bounded below by
zero, guaranteeing convergence of the \emph{objective value}.
However, convergence of the \emph{iterates} $\{\mat{U}^{(n)}_t\}$ to a
stationary point is not guaranteed in general due to non-convexity
\cite{de2000best}.

\paragraph{Computational cost.}
The dominant cost per HOOI iteration is $N$ truncated SVDs of matrices of
size $I_n \times \prod_{k \neq n} r_k$, giving a per-iteration complexity
of $O\!\left(N \cdot I \cdot r^{N-1}\right)$ for a balanced tensor with
$I_n = I$ and $r_n = r$.

% ------------------------------------------------------------
\subsection{Iteratively Reweighted Least Squares (IRLS) for $\ell_1$-Tucker}
\label{subsec:irls}
% ------------------------------------------------------------

\paragraph{Motivation.}
Direct minimization of $\mathcal{L}_1$ via subgradient methods converges
at rate $O(1/\sqrt{t})$, which is slow in practice.
IRLS sidesteps non-smoothness by solving a sequence of \emph{weighted}
Frobenius problems, each of which admits a standard ALS step.

\paragraph{IRLS reweighting.}
Let $r_{i_1 \cdots i_N} = x_{i_1 \cdots i_N} - \hat{x}_{i_1 \cdots i_N}$
be the current residual tensor.
Define element-wise weights
\begin{equation}
  w_{i_1 \cdots i_N}^{(t)}
  \;=\; \frac{1}{\bigl|r_{i_1 \cdots i_N}^{(t)}\bigr| + \varepsilon},
  \quad \varepsilon > 0 \text{ (smoothing parameter)},
  \label{eq:irls_weights}
\end{equation}
and the \emph{weighted} Frobenius problem
\begin{equation}
  \min_{\tensor{G},\, \{\mat{U}^{(n)}\}} \;
  \sum_{i_1, \ldots, i_N}
  w_{i_1 \cdots i_N}^{(t)}
  \bigl(x_{i_1 \cdots i_N} - \hat{x}_{i_1 \cdots i_N}\bigr)^2.
  \label{eq:irls_weighted}
\end{equation}
The connection to $\ell_1$ follows from the identity
\begin{equation}
  |r| = w \cdot r^2
  \quad \text{when } w = 1/|r|,
  \label{eq:irls_identity}
\end{equation}
which shows that minimizing~\eqref{eq:irls_weighted} with fixed weights
is a majorization step for $\mathcal{L}_1$.

\paragraph{Weighted ALS update.}
Let $\tensor{W}^{(t)}$ be the weight tensor and define the
element-wise product $\tilde{\tensor{X}} = \tensor{W}^{(t)} \odot \tensor{X}$.
The weighted mode-$n$ subproblem is
\begin{equation}
  \min_{\mat{U}^{(n)}} \;
  \bigl\|\tilde{\mat{X}}_{(n)} - \mat{U}^{(n)} \tilde{\mat{Z}}^{(n)}\bigr\|_F^2,
  \label{eq:irls_subproblem}
\end{equation}
which is solved by~\eqref{eq:als_update} applied to the weighted
unfolding $\tilde{\mat{X}}_{(n)}$.

Algorithm~\ref{alg:irls} presents the full IRLS procedure.

\begin{algorithm}[t]
  \caption{IRLS for $\ell_1$-Tucker decomposition}
  \label{alg:irls}
  \KwIn{Tensor $\tensor{X}$, rank $(r_1, \ldots, r_N)$,
        outer iterations $T$, inner ALS steps $K$,
        smoothing $\varepsilon > 0$}
  \KwOut{Core $\tensor{G}$, factor matrices $\{\mat{U}^{(n)}\}$}

  \BlankLine
  $\tensor{G}, \{\mat{U}^{(n)}\} \leftarrow \mathrm{HOOI}(\tensor{X},
    \mathbf{r})$
    \tcp*{warm start from Frobenius solution}

  \For{$t \leftarrow 1$ \KwTo $T$}{
    $\hat{\tensor{X}} \leftarrow
      \tensor{G} \times_1 \mat{U}^{(1)} \cdots \times_N \mat{U}^{(N)}$\;
    $\tensor{R} \leftarrow \tensor{X} - \hat{\tensor{X}}$
    \tcp*{residual}
    $\tensor{W} \leftarrow 1\,/\,(|\tensor{R}| + \varepsilon)$
    \tcp*{element-wise weights~\eqref{eq:irls_weights}}

    \For{$k \leftarrow 1$ \KwTo $K$}{
      \tcp{One pass of weighted ALS}
      \For{$n \leftarrow 1$ \KwTo $N$}{
        $\mat{U}^{(n)} \leftarrow$
          weighted ALS update~\eqref{eq:irls_subproblem}\;
      }
      $\tensor{G} \leftarrow
        \tensor{X} \times_1 {\mat{U}^{(1)}}^\top \cdots
                   \times_N {\mat{U}^{(N)}}^\top$\;
    }
  }
  \Return $\tensor{G},\; \{\mat{U}^{(n)}\}$\;
\end{algorithm}

\paragraph{Choice of $\varepsilon$.}
The parameter $\varepsilon$ prevents division by zero and controls the
approximation of $|\cdot|$ near the origin.
Smaller $\varepsilon$ gives a closer approximation to the true $\ell_1$
loss but may cause numerical instability.
A practical schedule starts with $\varepsilon = 10^{-1}$ and halves it
every few outer iterations.

% ------------------------------------------------------------
\subsection{Gradient-Based Optimization for Huber-Tucker}
\label{subsec:gradient}
% ------------------------------------------------------------

\paragraph{Motivation.}
The Huber loss $\mathcal{L}_H$ is everywhere differentiable, enabling
direct application of first- and second-order gradient methods.
We use this objective to compare gradient-based optimization with
ALS/IRLS in terms of convergence speed and final reconstruction quality.

\paragraph{Gradient derivation.}
Let $\tensor{R} = \tensor{X} - \hat{\tensor{X}}$ be the residual and
define the element-wise Huber pseudo-residual
\begin{equation}
  \psi_\delta(r)
  \;=\; \frac{\partial \rho_\delta(r)}{\partial r}
  \;=\; \begin{cases}
          r            & |r| \leq \delta, \\
          \delta\,\sign(r) & |r| > \delta.
        \end{cases}
  \label{eq:huber_grad}
\end{equation}
The gradient of $\mathcal{L}_H$ with respect to $\hat{\tensor{X}}$ is
\begin{equation}
  \frac{\partial \mathcal{L}_H}{\partial \hat{\tensor{X}}}
  \;=\; -\bm{\Psi}_\delta(\tensor{R}),
  \label{eq:loss_grad}
\end{equation}
where $\bm{\Psi}_\delta$ is applied element-wise.
By the chain rule and the mode-$n$ product structure,
\begin{equation}
  \frac{\partial \mathcal{L}_H}{\partial \mat{U}^{(n)}}
  \;=\; -\bm{\Psi}_\delta(\tensor{R})_{(n)}\,
         \left(
           \mat{U}^{(N)} \kron \cdots \kron \widehat{\mat{U}^{(n)}}
           \kron \cdots \kron \mat{U}^{(1)}
         \right)
         \mat{G}_{(n)}^\top,
  \label{eq:grad_U}
\end{equation}
where $\widehat{\cdot}$ denotes omission of the $n$-th factor.
The gradient with respect to the core is
\begin{equation}
  \frac{\partial \mathcal{L}_H}{\partial \mat{G}_{(n)}}
  \;=\; -{\mat{U}^{(n)}}^\top\,
         \bm{\Psi}_\delta(\tensor{R})_{(n)}\,
         \left(
           \mat{U}^{(N)} \kron \cdots \kron \mat{U}^{(1)}
         \right).
  \label{eq:grad_G}
\end{equation}

\paragraph{Optimization algorithm.}
We employ Adam~\cite{kingma2015adam} with parameters
$\alpha = 10^{-3}$, $\beta_1 = 0.9$, $\beta_2 = 0.999$.
All factor matrices are initialized from HOSVD to ensure a fair comparison
with HOOI.
The Huber threshold is set to $\delta = 0.1 \cdot \max|\tensor{R}_0|$,
where $\tensor{R}_0$ is the initial residual, so that $\delta$ adapts
to the scale of each tensor.

\paragraph{Implementation note.}
Equations~\eqref{eq:grad_U}--\eqref{eq:grad_G} can be evaluated without
autograd by computing mode-$n$ products explicitly, which avoids the
memory overhead of storing a computation graph over the full tensor.
For moderate-sized tensors (up to $256 \times 256 \times 3$), this
manual gradient implementation is preferred.

% ------------------------------------------------------------
\subsection{Comparison of Methods}
\label{subsec:opt_comparison}
% ------------------------------------------------------------

Table~\ref{tab:opt_comparison} summarizes the three optimization methods
along the axes most relevant to this project.

\begin{table}[h]
  \centering
  \caption{Comparison of optimization methods for Tucker approximation.}
  \label{tab:opt_comparison}
  \begin{tabular}{lcccc}
    \toprule
    Method & Objective & Convergence & Per-iter.\ cost & Hyperparams \\
    \midrule
    ALS / HOOI
      & $\mathcal{L}_F$
      & Monotone $\downarrow$, local
      & $O(N I r^{N-1})$
      & none \\
    IRLS
      & $\mathcal{L}_1$
      & Monotone (outer), local
      & $O(N I r^{N-1})$
      & $\varepsilon$, $K$ \\
    Adam (gradient)
      & $\mathcal{L}_H$
      & Non-monotone, local
      & $O(N I r^{N-1})$
      & $\alpha, \delta, \beta_{1,2}$ \\
    \bottomrule
  \end{tabular}
\end{table}

      \paragraph{Non-convexity and local minima.}
        All three methods are susceptible to local minima because the feasible set $\mathcal{M}_r$ is non-convex. To assess sensitivity to initialization, Section~\ref{sec:experiments} repeats each experiment with five random orthogonal initializations and reports the distribution of final reconstruction errors.\\

      \paragraph{Numerical stability of ALS.}
        The normal equation in~\eqref{eq:als_update} requires inverting $\mat{Z}^{(n)} \mat{Z}^{(n)\top}$. When this matrix is ill-conditioned --- which occurs when factor matrices are nearly linearly dependent --- a Tikhonov regularization term $\mu \mat{I}$ is added:

        \begin{equation}
          \mat{U}^{(n)} \leftarrow
          \unfold{X}{n} \mat{Z}^{(n)\top} \left(\mat{Z}^{(n)} \mat{Z}^{(n)\top} + \mu \mat{I}\right)^{-1}, \quad \mu = 10^{-8}
        \end{equation}

        In our experiments, the orthonormal constraint in HOOI makes this regularization unnecessary.\\

% % ============================================================
% %  sections/chapter_03_ko.tex
% %  3장: 최적화 방법
% % ============================================================

% \section{최적화 방법}
% \label{sec:optimization}

% Tucker 근사 문제~\eqref{eq:general_opt}는 인수 행렬 $\{\mat{U}^{(n)}\}$와
% 코어 텐서 $\tensor{G}$에 대해 \emph{결합 비볼록(jointly non-convex)}이다.
% 본 장에서는 이 비볼록성을 서로 다른 방식으로 다루는 세 가지 알고리즘군을
% 소개한다: 프로베니우스 목적 함수를 위한 교대 최소제곱법(ALS) 및 그
% 고차 정제인 HOOI, $\ell_1$ 목적 함수를 위한 반복 재가중 최소제곱법(IRLS),
% 그리고 Huber 목적 함수를 위한 경사 기반 방법.

% % ------------------------------------------------------------
% \subsection{교대 최소제곱법과 HOOI}
% \label{subsec:als}
% % ------------------------------------------------------------

% \paragraph{ALS 원리.}
% $\mathcal{L}_F$는 모든 인수를 동시에 고려하면 결합 비볼록이지만,
% 하나의 인수 행렬만 고정을 풀고 나머지를 고정하면
% \emph{이차(quadratic)} 함수 --- 즉 볼록 함수 --- 가 된다.
% ALS는 이 구조를 활용하여 인수를 순환하면서 각각의 최소제곱 문제를
% 정확하게 풀어나간다.

% \paragraph{모드-$n$ 부분 문제.}
% $\mat{U}^{(n)}$을 제외한 모든 인수 행렬을 고정하고 다음을 정의한다:
% \begin{equation}
%   \mat{Z}^{(n)}
%   \;\coloneqq\;
%   \mat{G}_{(n)}
%   \!\left(
%     \mat{U}^{(N)} \kron \cdots \kron \mat{U}^{(n+1)}
%     \kron \mat{U}^{(n-1)} \kron \cdots \kron \mat{U}^{(1)}
%   \right)^{\!\top}
%   \;\in\; \R^{r_n \times \prod_{k\neq n} I_k}.
%   \label{eq:Z_factor}
% \end{equation}
% 목적 함수는 다음으로 환원된다:
% \begin{equation}
%   \min_{\mat{U}^{(n)}} \;
%   \bigl\|\unfold{X}{n} - \mat{U}^{(n)} \mat{Z}^{(n)}\bigr\|_F^2,
%   \label{eq:als_subproblem}
% \end{equation}
% 이의 닫힌 형태 해는
% \begin{equation}
%   \mat{U}^{(n)}
%   \;\leftarrow\;
%   \unfold{X}{n} \mat{Z}^{(n)\top}
%   \!\left(\mat{Z}^{(n)} \mat{Z}^{(n)\top}\right)^{-1}.
%   \label{eq:als_update}
% \end{equation}
% $\mat{U}^{(n)}$의 열이 정규직교로 추가 제약되는 Tucker의 경우,
% 최적해는 $\unfold{X}{n} \mat{Z}^{(n)\top}$의 주요 $r_n$개
% 좌특이벡터이며, 이 경우 행렬 역연산이 불필요하다.

% \paragraph{HOOI 알고리즘.}
% 고차 직교 반복(HOOI)~\cite{de2000best}은 정규직교 제약 조건 하의
% ALS를 적용한다.
% 각 반복에서 수축 텐서(contracted tensor)
% \begin{equation}
%   \tensor{Y}^{(n)}
%   \;=\; \tensor{X}
%         \times_1 {\mat{U}^{(1)}}^\top
%         \cdots
%         \times_{n-1} {\mat{U}^{(n-1)}}^\top
%         \times_{n+1} {\mat{U}^{(n+1)}}^\top
%         \cdots
%         \times_N {\mat{U}^{(N)}}^\top
%   \label{eq:hooi_contracted}
% \end{equation}
% 를 구성하고(모드 $n$ 제외), 갱신은
% \begin{equation}
%   \mat{U}^{(n)}
%   \;\leftarrow\;
%   \mat{Y}^{(n)}_{(n)}\text{의 주요 }r_n\text{개 좌특이벡터}
%   \label{eq:hooi_update}
% \end{equation}
% 으로 이루어진다.
% 모든 인수 행렬이 수렴한 후 코어는 다음으로 복원된다:
% \begin{equation}
%   \tensor{G}
%   \;=\; \tensor{X}
%         \times_1 {\mat{U}^{(1)}}^\top
%         \times_2 {\mat{U}^{(2)}}^\top
%         \cdots
%         \times_N {\mat{U}^{(N)}}^\top.
%   \label{eq:hooi_core}
% \end{equation}

% 알고리즘~\ref{alg:hooi}에 전체 절차를 정리하였다.

% \begin{algorithm}[t]
%   \caption{Tucker 분해를 위한 HOOI (프로베니우스 노름)}
%   \label{alg:hooi}
%   \KwIn{텐서 $\tensor{X} \in \R^{I_1 \times \cdots \times I_N}$,
%         목표 랭크 $(r_1, \ldots, r_N)$,
%         허용 오차 $\varepsilon > 0$,
%         최대 반복 횟수 $T$}
%   \KwOut{코어 $\tensor{G}$, 인수 행렬 $\{\mat{U}^{(n)}\}$}

%   \BlankLine
%   \tcp{HOSVD로 초기화}
%   \For{$n \leftarrow 1$ \KwTo $N$}{
%     $\mat{U}^{(n)} \leftarrow$ $\unfold{X}{n}$의 주요-$r_n$ 좌특이벡터\;
%   }
%   \BlankLine
%   \tcp{반복 정제}
%   \For{$t \leftarrow 1$ \KwTo $T$}{
%     \For{$n \leftarrow 1$ \KwTo $N$}{
%       $\tensor{Y} \leftarrow
%         \tensor{X} \times_{k \neq n} {\mat{U}^{(k)}}^\top$\;
%       $\mat{U}^{(n)} \leftarrow$
%         $\unfold{Y}{n}$의 주요-$r_n$ 좌특이벡터\;
%     }
%     $\tensor{G} \leftarrow
%       \tensor{X} \times_1 {\mat{U}^{(1)}}^\top \cdots
%                  \times_N {\mat{U}^{(N)}}^\top$\;
%     $e_t \leftarrow
%       \frobn{\tensor{X} - \tensor{G} \times_1 \mat{U}^{(1)}
%              \cdots \times_N \mat{U}^{(N)}}$\;
%     \lIf{$|e_t - e_{t-1}| < \varepsilon$}{\textbf{종료}}
%   }
%   \Return $\tensor{G},\; \{\mat{U}^{(n)}\}$\;
% \end{algorithm}

% \paragraph{수렴성.}
% HOOI의 각 반복은 프로베니우스 재구성 오차를 단조 감소시킨다.
% 각 모드-$n$ 갱신이 나머지 인수를 고정한 상태에서
% 부분 문제~\eqref{eq:als_subproblem}의 정확한 최소화기이기 때문이다.
% 따라서 수열 $\{e_t\}$는 단조 비증가이며 0으로 하한이 있어
% \emph{목적 함수값}의 수렴이 보장된다.
% 그러나 비볼록성으로 인해 \emph{반복점} $\{\mat{U}^{(n)}_t\}$의
% 정류점(stationary point) 수렴은 일반적으로 보장되지 않는다~\cite{de2000best}.

% \paragraph{계산 복잡도.}
% HOOI 반복당 지배적인 비용은 크기 $I_n \times \prod_{k \neq n} r_k$인
% 행렬에 대한 $N$번의 절단 SVD이다.
% $I_n = I$, $r_n = r$인 균형 텐서에서 반복당 복잡도는
% $O\!\left(N \cdot I \cdot r^{N-1}\right)$이다.

% % ------------------------------------------------------------
% \subsection{$\ell_1$-Tucker를 위한 반복 재가중 최소제곱법 (IRLS)}
% \label{subsec:irls}
% % ------------------------------------------------------------

% \paragraph{동기.}
% 부분기울기(subgradient) 방법을 통한 $\mathcal{L}_1$의 직접 최소화는
% $O(1/\sqrt{t})$의 수렴률을 가지며 실제로 느리다.
% IRLS는 비매끄러움을 우회하기 위해 일련의 \emph{가중} 프로베니우스 문제를
% 반복적으로 풀며, 각각의 문제는 표준 ALS 단계로 해결된다.

% \paragraph{IRLS 재가중.}
% 현재 잔차 텐서를
% $r_{i_1 \cdots i_N} = x_{i_1 \cdots i_N} - \hat{x}_{i_1 \cdots i_N}$로
% 정의하고, 원소별 가중치를
% \begin{equation}
%   w_{i_1 \cdots i_N}^{(t)}
%   \;=\; \frac{1}{\bigl|r_{i_1 \cdots i_N}^{(t)}\bigr| + \varepsilon},
%   \quad \varepsilon > 0 \text{ (평활화 파라미터)}
%   \label{eq:irls_weights}
% \end{equation}
% 로 설정하면, \emph{가중} 프로베니우스 문제는
% \begin{equation}
%   \min_{\tensor{G},\, \{\mat{U}^{(n)}\}} \;
%   \sum_{i_1, \ldots, i_N}
%   w_{i_1 \cdots i_N}^{(t)}
%   \bigl(x_{i_1 \cdots i_N} - \hat{x}_{i_1 \cdots i_N}\bigr)^2
%   \label{eq:irls_weighted}
% \end{equation}
% 가 된다.
% $\ell_1$과의 연결은 항등식
% \begin{equation}
%   |r| = w \cdot r^2 \quad (w = 1/|r|\text{일 때})
%   \label{eq:irls_identity}
% \end{equation}
% 에서 비롯된다. 즉, 가중치를 고정하고~\eqref{eq:irls_weighted}를 최소화하는
% 것은 $\mathcal{L}_1$에 대한 우세화(majorization) 단계이다.

% \paragraph{가중 ALS 갱신.}
% $\tensor{W}^{(t)}$를 가중치 텐서, $\tilde{\tensor{X}} = \tensor{W}^{(t)} \odot \tensor{X}$를
% 원소별 곱으로 정의하면, 가중 모드-$n$ 부분 문제는
% \begin{equation}
%   \min_{\mat{U}^{(n)}} \;
%   \bigl\|\tilde{\mat{X}}_{(n)} - \mat{U}^{(n)} \tilde{\mat{Z}}^{(n)}\bigr\|_F^2
%   \label{eq:irls_subproblem}
% \end{equation}
% 이며, 식~\eqref{eq:als_update}를 가중 전개 $\tilde{\mat{X}}_{(n)}$에
% 적용하여 풀 수 있다.

% 알고리즘~\ref{alg:irls}에 전체 IRLS 절차를 정리하였다.

% \begin{algorithm}[t]
%   \caption{$\ell_1$-Tucker 분해를 위한 IRLS}
%   \label{alg:irls}
%   \KwIn{텐서 $\tensor{X}$, 랭크 $(r_1, \ldots, r_N)$,
%         외부 반복 $T$, 내부 ALS 단계 $K$,
%         평활화 $\varepsilon > 0$}
%   \KwOut{코어 $\tensor{G}$, 인수 행렬 $\{\mat{U}^{(n)}\}$}

%   \BlankLine
%   $\tensor{G}, \{\mat{U}^{(n)}\} \leftarrow \mathrm{HOOI}(\tensor{X}, \mathbf{r})$
%     \tcp*{프로베니우스 해로 웜 스타트}

%   \For{$t \leftarrow 1$ \KwTo $T$}{
%     $\hat{\tensor{X}} \leftarrow
%       \tensor{G} \times_1 \mat{U}^{(1)} \cdots \times_N \mat{U}^{(N)}$\;
%     $\tensor{R} \leftarrow \tensor{X} - \hat{\tensor{X}}$
%       \tcp*{잔차}
%     $\tensor{W} \leftarrow 1\,/\,(|\tensor{R}| + \varepsilon)$
%       \tcp*{원소별 가중치~\eqref{eq:irls_weights}}

%     \For{$k \leftarrow 1$ \KwTo $K$}{
%       \tcp{가중 ALS 1 패스}
%       \For{$n \leftarrow 1$ \KwTo $N$}{
%         $\mat{U}^{(n)} \leftarrow$
%           가중 ALS 갱신~\eqref{eq:irls_subproblem}\;
%       }
%       $\tensor{G} \leftarrow
%         \tensor{X} \times_1 {\mat{U}^{(1)}}^\top \cdots
%                    \times_N {\mat{U}^{(N)}}^\top$\;
%     }
%   }
%   \Return $\tensor{G},\; \{\mat{U}^{(n)}\}$\;
% \end{algorithm}

% \paragraph{$\varepsilon$ 선택.}
% $\varepsilon$은 0으로의 나눗셈을 방지하고 원점 근방에서 $|\cdot|$의
% 근사 정도를 조절한다.
% $\varepsilon$이 작을수록 진정한 $\ell_1$ 손실에 가까워지지만
% 수치 불안정의 위험이 증가한다.
% 실용적인 스케줄로는 $\varepsilon = 10^{-1}$에서 시작하여
% 몇 번의 외부 반복마다 절반으로 줄이는 방법을 사용한다.

% % ------------------------------------------------------------
% \subsection{Huber-Tucker를 위한 경사 기반 최적화}
% \label{subsec:gradient}
% % ------------------------------------------------------------

% \paragraph{동기.}
% Huber 손실 $\mathcal{L}_H$는 모든 점에서 미분 가능하므로
% 1차 및 2차 경사 방법을 직접 적용할 수 있다.
% 이 목적 함수를 사용하여 ALS/IRLS와 경사 기반 최적화를
% 수렴 속도 및 최종 재구성 품질 측면에서 비교한다.

% \paragraph{경사 유도.}
% $\tensor{R} = \tensor{X} - \hat{\tensor{X}}$를 잔차라 하고
% 원소별 Huber 유사잔차를
% \begin{equation}
%   \psi_\delta(r)
%   \;=\; \frac{\partial \rho_\delta(r)}{\partial r}
%   \;=\; \begin{cases}
%           r                    & |r| \leq \delta\text{인 경우,} \\
%           \delta\,\sign(r) & |r| > \delta\text{인 경우}
%         \end{cases}
%   \label{eq:huber_grad}
% \end{equation}
% 로 정의한다.
% $\mathcal{L}_H$의 $\hat{\tensor{X}}$에 대한 경사는
% \begin{equation}
%   \frac{\partial \mathcal{L}_H}{\partial \hat{\tensor{X}}}
%   \;=\; -\bm{\Psi}_\delta(\tensor{R})
%   \label{eq:loss_grad}
% \end{equation}
% 이며, 연쇄 법칙과 모드-$n$ 곱 구조로부터
% \begin{equation}
%   \frac{\partial \mathcal{L}_H}{\partial \mat{U}^{(n)}}
%   \;=\; -\bm{\Psi}_\delta(\tensor{R})_{(n)}\,
%          \left(
%            \mat{U}^{(N)} \kron \cdots \kron \widehat{\mat{U}^{(n)}}
%            \kron \cdots \kron \mat{U}^{(1)}
%          \right)
%          \mat{G}_{(n)}^\top
%   \label{eq:grad_U}
% \end{equation}
% ($\widehat{\cdot}$는 $n$번째 인수의 생략을 의미),
% 코어에 대한 경사는
% \begin{equation}
%   \frac{\partial \mathcal{L}_H}{\partial \mat{G}_{(n)}}
%   \;=\; -{\mat{U}^{(n)}}^\top\,
%          \bm{\Psi}_\delta(\tensor{R})_{(n)}\,
%          \left(
%            \mat{U}^{(N)} \kron \cdots \kron \mat{U}^{(1)}
%          \right)
%   \label{eq:grad_G}
% \end{equation}
% 이다.

% \paragraph{최적화 알고리즘.}
% Adam~\cite{kingma2015adam} 옵티마이저를 사용하며,
% 파라미터는 $\alpha = 10^{-3}$, $\beta_1 = 0.9$, $\beta_2 = 0.999$로 설정한다.
% HOOI와의 공정한 비교를 위해 모든 인수 행렬은 HOSVD로 초기화한다.
% Huber 임계값은 $\delta = 0.1 \cdot \max|\tensor{R}_0|$로 설정하여
% 각 텐서의 스케일에 맞게 적응시킨다.

% \paragraph{구현 참고.}
% 식~\eqref{eq:grad_U}--\eqref{eq:grad_G}은 모드-$n$ 곱을 명시적으로
% 계산함으로써 자동 미분(autograd) 없이 평가할 수 있다.
% 이 방식은 전체 텐서에 대한 계산 그래프 저장에 따른 메모리 오버헤드를
% 피할 수 있어, 중간 크기 텐서($256 \times 256 \times 3$ 이하)에서
% 수동 경사 구현이 선호된다.

% % ------------------------------------------------------------
% \subsection{방법 비교}
% \label{subsec:opt_comparison}
% % ------------------------------------------------------------

% 표~\ref{tab:opt_comparison}은 세 가지 최적화 방법을
% 본 프로젝트에서 가장 관련 있는 축으로 비교한다.

% \begin{table}[h]
%   \centering
%   \caption{Tucker 근사를 위한 최적화 방법 비교.}
%   \label{tab:opt_comparison}
%   \begin{tabular}{lcccc}
%     \toprule
%     방법 & 목적 함수 & 수렴 특성 & 반복당 비용 & 하이퍼파라미터 \\
%     \midrule
%     ALS / HOOI
%       & $\mathcal{L}_F$
%       & 단조 $\downarrow$, 국소
%       & $O(N I r^{N-1})$
%       & 없음 \\
%     IRLS
%       & $\mathcal{L}_1$
%       & 단조 (외부), 국소
%       & $O(N I r^{N-1})$
%       & $\varepsilon$, $K$ \\
%     Adam (경사)
%       & $\mathcal{L}_H$
%       & 비단조, 국소
%       & $O(N I r^{N-1})$
%       & $\alpha, \delta, \beta_{1,2}$ \\
%     \bottomrule
%   \end{tabular}
% \end{table}

% \paragraph{비볼록성과 국소 최솟값.}
% 세 방법 모두 실행 가능 집합 $\mathcal{M}_r$의 비볼록성으로 인해
% 국소 최솟값에 수렴할 수 있다.
% 초기화 민감도를 평가하기 위해, \ref{sec:experiments}장에서는
% 각 실험을 5개의 무작위 직교 초기화로 반복하고
% 최종 재구성 오차의 분포를 보고한다.

% \paragraph{ALS의 수치 안정성.}
% 식~\eqref{eq:als_update}의 정규 방정식은
% $\mat{Z}^{(n)} \mat{Z}^{(n)\top}$의 역행렬을 요구한다.
% 인수 행렬이 거의 선형 종속에 가까워질 때 이 행렬이 불량 조건(ill-conditioned)이
% 될 수 있으므로, Tikhonov 정규화 항 $\mu \mat{I}$을 추가한다:
% \begin{equation}
%   \mat{U}^{(n)}
%   \;\leftarrow\;
%   \unfold{X}{n} \mat{Z}^{(n)\top}
%   \!\left(\mat{Z}^{(n)} \mat{Z}^{(n)\top} + \mu \mat{I}\right)^{-1},
%   \quad \mu = 10^{-8}.
%   \label{eq:als_tikhonov}
% \end{equation}
% 본 실험에서는 HOOI의 정규직교 제약 조건이 이 정규화를 불필요하게 만든다.